\chapter{Masa (4/10)}

Grammatical descriptions are available for four of the ten languages of the Masa branch of Chadic languages -- two from the North Masa group and two from the South Masa group.

%The North Masa and South Masa languages contrast in terms of directional extensions, as can be seen in Table \ref{tab:Masa}. The two North Masa languages, Musey and Masana, have ventive and itive extensions that could plausibly be cognate. In contrast, the South Masa language Herde has just a ventive extension that does not appear related, and Pévé does not have directional extensions. In Pévé, directional meaning can instead be expressed in a serial verb construction \citet[185--189]{shayGrammarPeve2020}. 

%\begin{table}[H]
%	\caption{Masa group directional morphology}
%	\label{tab:Masa}
%	\begin{tabular}{lllllll} % add l for every additional column or remove as necessary
%		\lsptoprule
%		language & VENT & ITIVE  & UP & DOWN & INTO & OUT  \\ %table header
%		\midrule
%		Musey (muse1242) & \textit{-áj} & \textit{-há} & & & &  \\	
%		Masana (masa1322) & \textit{-ey} & 	\textit{-as} & \textit{kúló} & \textit{gà} & &  \\	
%		Herde (herd1236) & \textit{fi} & & & & &  \\	
%		Pévé (peve1243) & & & & & &  \\	
%		\lspbottomrule
%	\end{tabular}
%\end{table}

%The directional verbs across Masa languages share similar forms to express ventive motion, and a pattern of coronal segments can be seen across itive motion verbs. Other pairs of likely cognates can also be identified in Table \ref{tab:MasaV}. The North Masa languages use the same verb forms for inward/outward motion, while in South Masa languages, inward and outward motion are lexicalized as separate verbs. 

%\begin{table}[H]
%	\caption{Masa group directional verbs}
%	\label{tab:MasaV}
%	\begin{tabular}{lllllll} % add l for every additional column or remove as necessary
%		\lsptoprule
%		language & VENT & ITIVE  & UP & DOWN & INTO & OUT  \\ %table header
%		\midrule
%		Musey (muse1242) & \textit{mbara} & \textit{iira, ɓira} & \textit{diŋga, fulla} & \textit{diŋga, fulla} & \textit{sella} & \textit{sella} \\	
%		Masana (masa1322) & \textit{ma} & \textit{n} & & & \textit{cuk} & \textit{cuk}  \\	
%		Herde (herd1236) & \textit{mbu} & \textit{ta} & \textit{ur}  & & \textit{ze, cok} & \textit{pa, hùt} \\	
%		Pévé (peve1243) & \textit{mbə̀} & \textit{də̀} & \textit{dám} & \textit{dàm} & \textit{zye} & \textit{pã̰}  \\	
%		\lspbottomrule
%	\end{tabular}
%\end{table}

%In regard to caused accompanied motion (CAM), the South Masa languages both share a strategy of expanding the meaning of motion verbs through the use of a prepositional phrase. This differs from Masana where a verb glossed `take' combined with directional extensions to express CAM. (There is insufficent data on this topic for Musey.) 

%The two South Masa languages pattern together in using the same verb form to express BUY, as seen in Table \ref{tab:MasaBUY}, however, they have different strategies for expressing SELL: one a different verb stem, and another using a causative form of BUY. The two North Masa languages have similar forms for their verb for BUY, and both modify this form to express SELL, but do so with different morphs. In Musey \textit{ki} is described as marking ``\textit{achèvement}'' or completion, and in Masana \textit{góyò} is described as ``\textit{reversif}'' or reversive. 

%\begin{table}[H]
%	\caption{Masa group BUY-SELL verbs}
%	\label{tab:MasaBUY}
%	\begin{tabular}{lll} % add l for every additional column or remove as necessary
%		\lsptoprule
%		language & BUY & SELL    \\ %table header
%		\midrule
%		Musey (muse1242) & \textit{gussa} &  \textit{gus ki} [``\textit{achèvement}''] \\	
%		Masana (masa1322) & \textit{gìyàw}  & \textit{gìyàw góyò} [``\textit{reversif}''] \\	
%		Herde (herd1236) & \textit{tsòɓò} & \textit{ndùnò} \\	
%		Pévé (peve1243) & \textit{tsoɓ} & \textit{gi tsoɓ} [\textsc{caus}]\\	
%		\lspbottomrule
%	\end{tabular}
%\end{table}

\section{North Masa (2/6)}

Two North Masa languages are discussed here: Musey and Masana. Description of the morphosyntax of Gizey (gize1234) is in progress but not yet available. No relevant descriptions were found for Ham (hama1241), Marba (marb1239) or Zumaya (zuma1239).

\paragraph*{Musey (muse1242)} \label{muse}

\subparagraph{Directional verbs}~

\begin{table}[H]
	\caption{Musey directional verbs} 
	\label{tab:museV}
	\fittable{
	\begin{tabular}{lll} % add l for every additional column or remove as necessary
		\lsptoprule
		DIR & verb & gloss   \\ %table header
		\midrule
		MOTION & \textit{kalla} & `\textit{venir} [come], \textit{entrer} [enter], \textit{aller} [go], \textit{sortir} [exit], \\ 
		{}     &                & \textit{partir} [leave], \textit{dépasser} [pass]' \\
		ITIVE & \textit{iira, ɓira} & `\textit{aller} [go]' \\
		VENT & \textit{mbara} & `\textit{venir}, \textit{arriver} [come, arrive]' \\
		UP/DOWN & \textit{diŋga, fulla} & `\textit{monter}, \textit{descendre} [go up, go down]' \\
		INTO/OUT & \textit{sella} & `\textit{entrer}, \textit{sortir}, \textit{verser} [enter, exit, pour]' \\
		\lspbottomrule
	\end{tabular}
	}
\end{table}

%In addition to the motion verbs in Table \ref{tab:museV}, \citet[17]{shryockMagaaraHadVun1998} also lists a general motion verb with no apparent directionality: \textit{kalla} `venir [come], entrer [enter], aller [go], sortir [exit], partir [leave], dépasser [pass]'. 

\newpage
\subparagraph{Directional extensions}~

\begin{table}[H]
	\caption{Musey directional morphology}
	\label{tab:muse}
	\begin{tabular}{llll} % add l for every additional column or remove as necessary
		\lsptoprule
		forms & source gloss & direction  & other functions  \\ %table header
		\midrule
		\textit{-áj}  &  \textsc{dir}  & \textsc{vent}  &   \\
		\textit{-há}  & \textsc{dir}  & \textsc{itive} &     \\	
		\lspbottomrule
	\end{tabular}
\end{table}

The description of Musey grammar by \citet{platielEsquisseEtudeMusey1968} does not include any mention of motion-related verbal morphology, but \citet{dassidiMorphophonologicalStructureMusej2015} describes a ventive suffix \textit{-áj}, as in (\ref{ex2:muse1}), and an itive suffix \textit{-há}, as in (\ref{ex2:muse2}). Only a few examples are given and they do not generally make explicit what the semantic contribution of the directional suffixes are. 

\ea\label{ex2:muse1}
\langinfo{Musey verb \textit{līŋ} `run' with ventive marker}{}{\cite[76]{dassidiMorphophonologicalStructureMusej2015}} \\
\gll Sánàmā līŋ-áj tá ànú	\\
Sanama run-\gl{vent} body 1\gl{sg}.\gl{obj} \\
\glt `Sanama ran towards me.'

\ex\label{ex2:muse2}
\langinfo{Musey verb \textit{kál} `go/come' with itive marker}{}{\cite[76]{dassidiMorphophonologicalStructureMusej2015}} \\
\gll ndàt kál-há bùdúwā	\\
3\gl{sg}.\gl{f} go/come-\gl{itive} outside		\\
\glt `She went out.'
\z 

\subparagraph{Caused accompanied motion (CAM)}

Several caused accompanied motion verbs are listed in the Musey lexicon, but it is unclear from the available sources what strategies are used to express directed CAM. 

\subparagraph{Buy-sell verbs}

The lexicon entry for the verb \textit{gussa} `buy' also lists a modified form \textit{gus ki} `sell' \citep[13]{shryockMagaaraHadVun1998}. The morpheme \textit{ki} is defined as `\textit{mot qui marque l'achèvement d'une action} [word which marks the completion of an action]' \citep[18]{shryockMagaaraHadVun1998}.

\newpage
\paragraph*{Masana (masa1322)}  \label{masa}

\subparagraph{Directional verbs}~

\begin{table}[H]
	\caption{Masana directional verbs} 
	\label{tab:masaV}
	\begin{tabular}{lll} % add l for every additional column or remove as necessary
		\lsptoprule
		DIR & verb & gloss   \\ %table header
		\midrule
		MOTION & \textit{n} & `\textit{aller} [go], \textit{venir} [come]' \\
		VENT & \textit{ma} & `\textit{venir} [come]' \\
		INTO/OUT & \textit{cuk} & `\textit{entrer} [enter], \textit{sortir} [exit]' \\
		\lspbottomrule
	\end{tabular}
\end{table}

\subparagraph{Directional extensions}~

\begin{table}[H]
	\caption{Masana directional morphology}
	\label{tab:masa}
	\begin{tabular}{llll} % add l for every additional column or remove as necessary
		\lsptoprule
		forms & source gloss & direction  & other functions  \\ %table header
		\midrule
		\textit{-ey}  &  $\rightarrow$O  & \textsc{vent}  &   \\
		\textit{-as}  &  ͏͏$\leftarrow$O  & \textsc{itive} &     \\
%		\textit{kúló}  &  ͏͏$\uparrow$  & \textsc{dir.up} &  away   \\		
%		\textit{gà}  &  ͏͏$\downarrow$  & \textsc{dir.down} &  away   \\		
		\lspbottomrule
	\end{tabular}
\end{table}

\citet{melisDescriptionMasaTchad1999} describes four directional particles. The ventive extension \textit{-ey} and the itive extension \textit{-as} show signs of being phonologically attached to the preceding word which is often the verb, but not always (see \ref{ex2:masa9}). They can have directional meaning with intransitive motion verbs, as in (\ref{ex2:masa1}) and (\ref{ex2:masa2}), and transitive motion verbs, as in (\ref{ex2:masa3}) and (\ref{ex2:masa4}).

\ea\label{ex2:masa1}
\langinfo{Masana verb \textit{n} `go/come' with ventive marker}{}{\cite[215]{melisDescriptionMasaTchad1999}} \\
\gll  n-àŋ-ìyà	\\
go/come-you-\gl{vent}	\\
\glt `come (toward me)' [French: `\textit{viens} (\textit{vers moi})']

\ex\label{ex2:masa2}
\langinfo{Masana verb \textit{n} `go/come' with itive marker}{}{\cite[216]{melisDescriptionMasaTchad1999}} \\
\gll n-àk-ás kàt ɓùs-ùk-kà	\\
go/come-you-\gl{itive} toward sister-you-\gl{def}	\\
\glt `Go away toward your sister (far from me).' \newline	[French: `\textit{Va-t-en auprès de ta soeur} (\textit{loin de moi}).']

\ex\label{ex2:masa3}
\langinfo{Masana verb \textit{vùl} `give' with ventive marker}{}{\cite[213]{melisDescriptionMasaTchad1999}} \\
\gll gór-t vùl-àn-t-áy m-àn mày	\\
small-\gl{def} give.\gl{ipfv}.1\gl{sg}-3\gl{sg}.\gl{f}-\gl{vent} \gl{prep}.1\gl{sg} and	\\
\glt `...and you give the small one to me.' \newline [French: `...\textit{et tu me donnes la petite pour moi}.']

\ex\label{ex2:masa4}
\langinfo{Masana verb \textit{vùl} `give' with itive marker}{}{\cite[214]{melisDescriptionMasaTchad1999}} \\
\gll nàʔ vúl-úm-t-àsà	\\
3\gl{sg}.\gl{f} give.\gl{pfv}-3\gl{sg}.\gl{m}-3\gl{sg}.\gl{f}-\gl{itive}	\\
\glt `She gives it to him (what is given leaves the subject).' \newline [French: `\textit{Elle le lui donna} (\textit{ce qui est donné quitte le sujet}).']
\z 

In addition, the ventive and itive extensions are said to have various aspectual functions mostly relating to viewing the activity as completed \citep[216--220]{melisDescriptionMasaTchad1999}. It is not always clear what context gives rise to various interpretations, even with motion verbs. In (\ref{ex2:masa5}) the ventive extension has a directional interpretation while the ventive extension with the same verb in (\ref{ex2:masa6}) is said to have an aspectual (non-ventive) interpretation. 

\ea\label{ex2:masa5}
\langinfo{Masana verb \textit{kàl} `leave' with ventive marker}{}{\cite[214]{melisDescriptionMasaTchad1999}} \\
\gll kàl-àŋ-ày jìyà	\\
leave-2\gl{sg}-\gl{vent} outside	\\
\glt `go out (come toward me)' [French: `\textit{sors} (\textit{viens vers moi})']

\ex\label{ex2:masa6}
\langinfo{Masana verb \textit{kàl} `leave' with ventive marker}{}{\cite[214]{melisDescriptionMasaTchad1999}} \\
\gll bún kàl-èy lúmúʔù	\\
father-2\gl{sg}.\gl{poss} leave.\gl{pfv}-\gl{vent} market	\\
\glt `My father has gone to the market.' \newline [French: `\textit{Mon père est parti au marché}.']
\z

There are forms identical to the itive and ventive extension (and with the same glosses) which are found in non-verbal clauses \citep[221]{melisDescriptionMasaTchad1999}. In these cases the morphemes are presented as syntactically independent words, and they do not appear to have any motion-related function. For these reasons, it appears that although the forms are superficially similar, they can be treated as two separate parts of the grammar; one an extension and one a preposition or particle.

In contrast, the upward directional extension \textit{kúló} and the downward directional extension \textit{gà} are always presented as separate words (rather than suffixes/clitics), as in (\ref{ex2:masa7}) and (\ref{ex2:masa8}). They can also appear in nonverbal clauses in which case they have similar (perhaps locational) meaning. Given that the differences in their use in verbal and nonverbal clauses are not particularly significant, these words are here considered as adverbs rather than as part of the verbal morphology.

\ea\label{ex2:masa7}
\langinfo{Masana verb \textit{ɦál} `jump' with upward marker}{}{\cite[226]{melisDescriptionMasaTchad1999}} \\
\gll ɮó ɦál kúló mìrìk mìrìk 	\\
ɮo jump.\gl{pfv} up \gl{ideo}	\\
\glt `ɮo hopped on both feet.' [French: `\textit{ɮo sauta à pieds joints}.']

\ex\label{ex2:masa8}
\langinfo{Masana verb \textit{ɦál} `jump' with downward marker}{}{\cite[226]{melisDescriptionMasaTchad1999}} \\
\gll  ʔàm ɦál gàʔà	 \\
3\gl{sg}.\gl{m} jump.\gl{pfv} down	\\
\glt `He jumped down.' [French: `\textit{Il descendit d’un bond}.']
\z 

The word \textit{góy} is also described as a directional particle with a meaning overlapping with the itive marker \textit{-as} \citep[223]{melisDescriptionMasaTchad1999}, but in most cases the examples can be treated as either relating to locative meaning (`far') or other non-motion meaning.

\subparagraph{Caused accompanied motion (CAM)}

In addition to general caused accompanied motion verbs like \textit{nàr} `carry' and \textit{màr} `carry', the ventive and itive extensions combine with various verbs glossed `take' to express directed CAM, as in (\ref{ex2:masa9}) and (\ref{ex2:masa10}). The translations of these examples are not clearly ventive. \citet[216]{melisDescriptionMasaTchad1999} explains this by saying that in this case the deictic center is the grammatical subject rather than the speaker. 

\ea\label{ex2:masa9}
\langinfo{Masana verb \textit{yòw} `take' with ventive marker}{}{\cite[216]{melisDescriptionMasaTchad1999}} \\
\gll ʔísí yòw díf-n-ìyà	\\
3\gl{pl} take.\gl{pfv} flute-\gl{def}-\gl{vent}	\\
\glt `They took the flutes with them.' [French: `\textit{Ils prirent les flûtes avec eux}.']

\ex\label{ex2:masa10}
\langinfo{Masana verb \textit{v} `take' with ventive marker}{}{\cite[223]{melisDescriptionMasaTchad1999}} \\
\gll v-àk-kà-s tùn-àk-kà góyò	\\
take-2\gl{sg}-3\gl{sg}-\gl{vent} put-2\gl{sg}-3\gl{sg} \textsc{reversive}	\\
\glt `Take it away far from me and put it elsewhere.' \newline [French: `\textit{Emporte-le loin de moi et pose-le ailleurs}.']
\z 	

\subparagraph{Buy-sell verbs}

\citet[162, cf. 176]{melisDictionnaireMasafrancaisDialectes2006} lists the verb \textit{gìyàw} `buy' along with the form \textit{gìyàw góyò} `sell'. The marker \textit{góy(o)} is labeled \textit{reversif}: ``\textit{le reversif s'organise par rapport à un seul repère qui correspond à la situation du terme source au moment de l'énonciation et va servir à indiquer un changement de cette situation. Ce changement de situation peut correspondre à un changement de localisation or d'état.} [the reversive is organized in relation to a single marker which corresponds to the situation of the source term at the time of the utterance and is used to indicate a change in this situation. This change of situation may correspond to a change of location or state]'' \citep[168]{melisDictionnaireMasafrancaisDialectes2006}. See also \citet[224]{melisDescriptionMasaTchad1999}.

\section{South Masa (2/4)}

Two South Masa languages, Herde and Pévé, are included here. No relevant publications were found for Ngete (nget1241) and the publications on Mesme (mesm1239) either do not cover the verbal morphology or were inaccessible.

\paragraph*{Herde (herd1236)} \label{herd}

\subparagraph{Directional verbs}~

\begin{table}[H]
	\caption{Herde directional verbs} 
	\label{tab:herdV}
	\begin{tabular}{lll} % add l for every additional column or remove as necessary
		\lsptoprule
		DIR & verb & gloss   \\ %table header
		\midrule
		ITIVE & \textit{ta} & `\textit{aller} [go]' \\
		VENT & \textit{mbu} & `\textit{venir} [come]' \\
		INTO& \textit{ze, cok} & `\textit{entrer} [enter]' \\
		OUT & \textit{pa, hùt} & `\textit{sortir} [exit]' \\
		UP & \textit{ur} & `\textit{monter} [go up]' \\
		\lspbottomrule
	\end{tabular}
\end{table}

Verbs in Table \ref{tab:herdV} are from \citet{sachnineLexiqueZimefrancaisVun1985}. 

\subparagraph{Directional extensions}~

\begin{table}[H]
	\caption{Herde directional morphology}
	\label{tab:herd}
	\begin{tabular}{llll} % add l for every additional column or remove as necessary
		\lsptoprule
		forms & source gloss & direction  & other functions  \\ %table header
		\midrule
		\textit{fi}  &  \textsc{direc.}  & \textsc{vent}  &   \\
		\lspbottomrule
	\end{tabular}
\end{table}

\citet[290]{sachnineLame1982} includes an entry for a directional marker \textit{fi} that ``\textit{indique un mouvement de retour vers le point de départ} [indicates a movement of return to the point of departure]''. In the few examples where this marker is found, it appears to express ventive direction, as in (\ref{ex2:herd1}) and (\ref{ex2:herd2}).

\ea\label{ex2:herd1}
\langinfo{Herde verb \textit{tū} `hurry' with ventive marker}{}{\cite[362]{sachnineLame1982}} \\
\gll  ńdā ɬè-tū ńdā fí ɬé ˀálá mbù zàˀán dìˀí nā ˀiá ɓō 	\\
she ??-hurry she \gl{vent} ?? to come find.me before I leave after \\
\glt `She hurried to come see before I left.' \newline [French: `\textit{elle s’est hâtée pour venir me voir avant que je ne parte}.']

\ex\label{ex2:herd2}
\langinfo{Herde verb \textit{dzìk} `push' with ventive marker}{}{\cite[408]{sachnineLame1982}} \\
\gll kə́ra?n ˀá mbə̀mbə̀ra dzìk fí mbù ˀá.ndì \\
basket with weight push \gl{vent} come with \\
\glt `The basket is very heavy, it will be necessary to drag it here.' \newline [French: `\textit{Le panier est très lourd, il faudra l ’amener en le traînant}.']
\z 

\subparagraph{Caused accompanied motion (CAM)}

The lexical entry for \textit{mbù} `\textit{venir} [come]' includes an entry for (presumably ventive) caused accompanied motion when the motion verb is followed by the preposition \textit{ˀá} `\textit{avec} [with]' meaning `\textit{apporter} [bring]' as in (\ref{ex2:herd3}).

\ea\label{ex2:herd3}
\langinfo{Herde verb \textit{mbù} `come' with prepositional phrase}{}{\cite[362]{sachnineLame1982}} \\
\gll  ńdā mbù-ˀá fūn nè ˀá vùtù \\
she come-with boule give to chief  \\
\glt `She brought the food to the chief.' \newline [French: `\textit{Elle a apporté de la nourriture au chef}.']
\z

\subparagraph{Buy-sell verbs}

Herde has two verbs: \textit{tsòɓò} `buy' and \textit{ndùnò} `sell' \citep[373, 401]{sachnineLame1982}. 

\newpage
\paragraph*{Pévé (peve1243)} \label{peve}

Pévé is described in a grammar by \citet{shayGrammarPeve2020}. 

\subparagraph{Directional verbs}~

\begin{table}[H]
	\caption{Pévé directional verbs} 
	\label{tab:peveV}
	\begin{tabular}{lll} % add l for every additional column or remove as necessary
		\lsptoprule
		DIR & verb & gloss   \\ %table header
		\midrule
		ITIVE & \textit{də̀} & `go' \\
		VENT & \textit{mbə̀} & `come' \\
		UP & \textit{dám} & `ascend' \\
		DOWN & \textit{dàm} & `descend' \\
		INTO & \textit{zye} & `enter' \\
		OUT & \textit{pã̰} & `exit' \\
		\lspbottomrule
	\end{tabular}
\end{table}

The verb \textit{dám} is defined as `to ascend, rise, climb out of' \citep[333, 350]{shayGrammarPeve2020} and is used in one example sentence \citep[145]{shayGrammarPeve2020}. The verb \textit{dàm} (with a low tone) is defined as `descend, climb out of' \citep[346]{shayGrammarPeve2020}, but does not appear in any example sentences. 

\subparagraph{Directional extensions}~

Pévé does not have any verbal morphology expressing directional meaning.

\subparagraph{Caused accompanied motion (CAM)}

Ventive caused accompanied motion can be expressed by the verb \textit{mbə́} `come' followed by an associative prepositional phrase, as in (\ref{ex2:peve7}). It is not known if the same strategy can be employed for itive caused accompanied motion.

\ea\label{ex2:peve7}
\langinfo{Pévé ventive CAM with associative preposition}{}{\cite[217]{shayGrammarPeve2020}} \\
\gll ha mbə́ kə sə̀lày si su \\
2\gl{m} come \gl{asoc} money \textsc{assertive} \gl{q} \\
\glt `Did you bring (lit. ‘come with’) money?'
\z 

%\ea\label{ex2:peve8}
%\langinfo{Pévé ventive CAM in SVC}{}{\cite[244]{shayGrammarPeve2020}} \\
%\gll swə mə̀ ni ɬé nàn mbə́ ti kə ndi nə̀ bay màn-à \\
%man REL PRO take 1SG.O come here ASSC OP POST friend 1SG.POSS-FV \\
%\glt `the man who brought me here is my friend.'
%\z

\subparagraph{Buy-sell verbs}

The verb \textit{tsoɓ} means `buy' and there is no dedicated lexical item for `sell'. Rather, a periphrastic causative \textit{gi tsoɓ} `make buy' (also referred to as a ``compound verb'') is used to express SELL, as in (\ref{ex2:peve8}) \citep[70, 294]{shayGrammarPeve2020}.

\ea\label{ex2:peve8}
\langinfo{Pévé multiword expression for `sell'}{}{\cite[59]{shayGrammarPeve2020}} \\
\gll  na ráʔ tímbì kunə kə́dàn gi tsoɓ rum(-u) \\
1\gl{sg} take calabash \gl{dem}.\gl{pl} \gl{purp} make purchase 3\gl{sg}(-\textsc{final.vowel}) \\
\glt `I took the calabashes in order to sell them.'
\z
